\IntroVorlesung
\writeSectionFrame{\presentationTitle}{}

% Generelle Phasen des Datenbankentwurfs
\begin{frame}{Phasen des Datenbankentwurfs}
    \begin{enumerate}    
        \item<1-> Anforderungsanalyse: Informale Beschreibung 
        \item<2-> Konzeptueller Entwurf: Entity-Relationship-Model 
        \item<3-> Implementationsentwurf:
              \begin{itemize}    
                  \item Relationales Schema
                  \item Normalisierung
              \end{itemize}
        \item<4-> Physischer Entwurf:
              \begin{itemize}    
                  \item Indexstrukturen
                  \item Technische Einstellung optimieren
              \end{itemize}
    \end{enumerate}
    \begin{center}    
        \begin{tikzpicture}    
            % first image
            \onslide<1->{    
                    \node[anchor=east] at ([xshift=-0.1cm]0,0) {    
                        
\includegraphics[height=1.1cm]{\BILDERPFAD/check-list.png}
                    };
                }
            
            % first arrow
            \onslide<2->{    
                    \draw[    
                        draw=Blue,
                        line width=2mm,
                            -{Triangle[length=3mm,width=3mm]}
                    ] (0,0) -- (0.7cm,0);
                }
            
            % second image
            \onslide<2->{    
                    \node[anchor=west] at (0.7cm,0){    
                        \includegraphics[height=1.7cm]{\BILDERPFAD/ERM_Example.png}
                    };
                }
            
            % second arrow
            \onslide<3->{    
                    \draw[    
                        draw=Blue,
                        line width=2mm,
                            -{Triangle[length=3mm,width=3mm]}
                    ] (3.3cm,0) -- (4.0cm,0);
                }
            
            % third image
            \onslide<3->{    
                    \node[anchor=west] at (4.1cm,0){    
                        \includegraphics[height=1.4cm]{\BILDERPFAD/Relationales_Schema_Example.png}
                    };
                }
            
            % third arrow
            \onslide<4->{    
                    \draw[    
                        draw=Blue,
                        line width=2mm,
                            -{Triangle[length=3mm,width=3mm]}
                    ] (7.3cm,0) -- (8.0cm,0);
                }
            
            % fourth image
            \onslide<4->{    
                    \node[anchor=west] at (8.1cm,0){    
                        \includegraphics[height=1.0cm]{\BILDERPFAD/parameter_tuning.png}
                    };
                }
        \end{tikzpicture}
    \end{center}
\end{frame}

% Entity-Relationship-Modell
\begin{frame}{Entity-Relationship-Modell}
    
    \begin{tikzpicture}[    
            scale=0.68,
            transform shape,
            entity/.style={    
                    draw,
                    rectangle,
                    minimum width=2.4cm,
                    minimum height=0.9cm,
                    align=center,
                    font=\scriptsize
                },
            relation/.style={    
                    draw,
                    diamond,
                    aspect=2,
                    minimum width=2.1cm,
                    minimum height=0.9cm,
                    align=center,
                    font=\scriptsize
                },
            attribute/.style={    
                    draw,
                    ellipse,
                    minimum width=1.6cm,
                    minimum height=0.7cm,
                    align=center,
                    font=\scriptsize
                },
            line/.style={draw, -}
        ]
        
        % Entitäten
        \node[entity] (studenten) at (0,4) {Studenten};
        \node[entity] (vorlesungen) at (8,4) {Vorlesungen};
        \node[entity] (assistenten) at (0,-2) {Assistenten};
        \node[entity] (professoren) at (8,-2) {Professoren};
        
        % Beziehungen
        \node[relation] (hoeren) at (4,4) {hören};
        \node[relation] (pruefen) at (4,2) {prüfen};
        \node[relation] (lesen) at (8,1.5) {lesen};
        \node[relation] (arbeitenfuer) at (4,-2) {arbeiten für};
        
        % Attribute Studenten
        \node[attribute] (smatr) at (-2.1,5.2) {\underline{MatrNr}};
        \node[attribute] (svorname) at (-2.6,4.3) {Vorname};
        \node[attribute] (sname) at (-2.6,3.3) {Name};
        \node[attribute] (ssemester) at (-2.1,2.2) {Semester};
        
        % Attribute Vorlesungen
        \node[attribute] (vnr) at (10.0,5.2) {\underline{VorlNr}};
        \node[attribute] (vsws) at (10.2,4.3) {SWS};
        \node[attribute] (vtitel) at (10.0,3.3) {Titel};
        
        % Attribute Assistenten
        \node[attribute] (apers) at (-2.2,-0.9) {\underline{PersNr}};
        \node[attribute] (avorname) at (-3.0,-1.8) {Vorname};
        \node[attribute] (aname) at (-3.0,-2.7) {Name};
        \node[attribute] (ataetigkeit) at (-2.3,-3.8) {Tätigkeit};
        
        % Attribute Professoren
        \node[attribute] (ppers) at (6.6,-3.8) {\underline{PersNr}};
        \node[attribute] (pvorname) at (8.2,-4.2) {Vorname};
        \node[attribute] (pfach) at (10.1,-3.9) {Fach};
        \node[attribute] (pname) at (11.2,-2.5) {Name};
        
        % Attribut Beziehung prüfen
        \node[attribute] (note) at (4,0.5) {Note};
        
        % Kanten Beziehungen
        \draw[line] (studenten) -- (hoeren);
        \draw[line] (hoeren) -- (vorlesungen);
        
        \draw[line] (studenten) -- (pruefen);
        \draw[line] (pruefen) -- (vorlesungen);
        \draw[line] (pruefen) -- (professoren);
        
        \draw[line] (vorlesungen) -- (lesen);
        \draw[line] (lesen) -- (professoren);
        
        \draw[line] (assistenten) -- (arbeitenfuer);
        \draw[line] (arbeitenfuer) -- (professoren);
        
        % Kanten Attribute
        \draw[line] (studenten) -- (smatr);
        \draw[line] (studenten) -- (svorname);
        \draw[line] (studenten) -- (sname);
        \draw[line] (studenten) -- (ssemester);
        
        \draw[line] (vorlesungen) -- (vnr);
        \draw[line] (vorlesungen) -- (vsws);
        \draw[line] (vorlesungen) -- (vtitel);
        
        \draw[line] (assistenten) -- (apers);
        \draw[line] (assistenten) -- (avorname);
        \draw[line] (assistenten) -- (aname);
        \draw[line] (assistenten) -- (ataetigkeit);
        
        \draw[line] (professoren) -- (ppers);
        \draw[line] (professoren) -- (pvorname);
        \draw[line] (professoren) -- (pname);
        \draw[line] (professoren) -- (pfach);
        
        \draw[line] (pruefen) -- (note);
        
    \end{tikzpicture}
    
\end{frame}

%%%%%%%%%%%%%%%%%%%%%%%%%%%%%%%%%%%%%%%%%%%%%%%%%%%%%%%%%%%%%%%%%%%%%%%%%%%%%%%%
% Entity-Typen
\begin{frame}{Entity-Typen}
    \begin{columns}[c]
        \column{0.55\textwidth}
        \begin{itemize}    
            \item<2-> Entitäten: Gegenstände, die wir in DB speichern wollen.
            \item<3-> Entity-Typen: Fassen Entitäten zusammen.
            \item<4-> Attribute: beschreibende Werte der Entitäten
            \item<5-> Schlüssel: Attribut(e), die die Entität eindeutig bestimmen.
            \item<6-> immer plural
        \end{itemize}
        
        \column{0.45\textwidth}
        \centering
        \only<1->{    
                \begin{tikzpicture}[    
                        entity/.style={    
                                draw,
                                rectangle,
                                minimum width=2.4cm,
                                minimum height=0.9cm,
                                align=center,
                                font=\scriptsize
                            },
                        attribute/.style={    
                                draw,
                                ellipse,
                                minimum width=1.6cm,
                                minimum height=0.7cm,
                                align=center,
                                font=\scriptsize
                            },
                        line/.style={draw, -}
                    ]
                    
                    % Entität
                    \node[entity] (vorlesungen) at (0,0) {Vorlesungen};
                    
                    % Attribute
                    \node[attribute] (vnr) at (0,1.8) {\underline{VorlNr}};
                    \node[attribute] (vsws) at (2.8,0) {SWS};
                    \node[attribute] (vtitel) at (0,-1.8) {Titel};
                    
                    % Kanten
                    \draw[line] (vorlesungen) -- (vnr);
                    \draw[line] (vorlesungen) -- (vsws);
                    \draw[line] (vorlesungen) -- (vtitel);
                    
                \end{tikzpicture}
            }
        
    \end{columns}
\end{frame}

%%%%%%%%%%%%%%%%%%%%%%%%%%%%%%%%%%%%%%%%%%%%%%
% Entität vs Entity-Typ

\begin{frame}{Entität vs Entity-Typ}
    
\end{frame}


%%%%%%%%%%%%%%%%%%%%%%%%%%%%%%%%%%%%%%%%%%%%%%
% Beziehungstypen

\begin{frame}{Beziehungstypen}
    
    \begin{itemize}    
        \item<2-> Beziehungstypen beschreiben Zusammenhänge zwischen Entity-Typen
        \item<3-> Beispiel: Professoren \textbf{lesen} Vorlesungen
        \item<4-> Mathematisch:\\
              $\text{lesen} \subseteq \text{Professoren} \times \text{Vorlesungen}$
        \item<5-> Allgemein für $n$ Entitäten:\\
              $\text{R} \subseteq E_1 \times E_2 \times \cdots E_n$
    \end{itemize}
    
    \vspace{0.3cm}
    
    \centering
    \only<1->{    
            \begin{tikzpicture}    
                
                % Entitäten
                \node[entity] (prof) at (0,0) {Professoren};
                \node[entity] (vl) at (6,0) {Vorlesungen};
                
                % Beziehung
                \node[relation] (lesen) at (3,0) {lesen};
                
                % Attribute Professoren
                \node[attribute] (persnr) at (0,1.3) {\underline{PersNr}};
                \node[attribute] (pname) at (0,-1.3) {Name};
                
                % Attribute Vorlesungen
                \node[attribute] (vorlnr) at (6,1.3) {\underline{VorlNr}};
                \node[attribute] (raum) at (6,-1.3) {Raum};
                
                % Kanten
                \draw (prof) -- (lesen);
                \draw (lesen) -- (vl);
                
                \draw (prof) -- (persnr);
                \draw (prof) -- (pname);
                
                \draw (vl) -- (vorlnr);
                \draw (vl) -- (raum);
                
            \end{tikzpicture}
        }
\end{frame}

%%%%%%%%%%%%%%%%%%%%%%%%%%%%%%%%%%%%%%%%%%%%%%%%%%%%%%%%%%%%
% Kardinalitäten: 1:1
\begin{frame}{Kardinalitäten: 1 zu 1}
    \centering
    \begin{tikzpicture}    
        
        % Entitäten
        \node[entity] (prof) at (0,0) {$E_1$};
        \node[entity] (vl) at (6,0) {$E_2$};
        
        % Beziehung
        \node[relation] (lesen) at (3,0) {$R$};
        
        % Kardinalitäten
        \node at (1.3,0.2) {\small 1};
        \node at (4.7,0.2) {\small 1};
        
        % Kanten
        \draw (prof) -- (lesen);
        \draw (lesen) -- (vl);
        
    \end{tikzpicture}
    \vspace{0.5cm}
    \begin{itemize}    
        \item Jeder Entität aus $E_1$ ist höchstens ein $E_2$ zugeordnet.
        \item Jeder Entität aus $E_2$ ist höchstens ein $E_1$ zugeordnet.
    \end{itemize}
    \vspace{0.5cm}
    \centering
    \begin{tikzpicture}    
        
        % Entitäten
        \node[entity] (prof) at (0,0) {Person};
        \node[entity] (vl) at (6,0) {Reisepass};
        
        % Beziehung
        \node[relation] (lesen) at (3,0) {besitzt};
        
        % Kardinalitäten
        \node at (1.3,0.2) {\small 1};
        \node at (4.7,0.2) {\small 1};
        
        % Kanten
        \draw (prof) -- (lesen);
        \draw (lesen) -- (vl);
    \end{tikzpicture}
    
\end{frame}

%%%%%%%%%%%%%%%%%%%%%%%%%%%%%%%%%%%%%%%%%%%%%%%%%%%%%%%%%%%%
% Kardinalitäten: 1:N
\begin{frame}{Kardinalitäten: 1 zu $N$}
    
    \centering
    
    \begin{tikzpicture}    
        
        % Entitäten
        \node[entity] (prof) at (0,0) {$E_1$};
        \node[entity] (vl) at (6,0) {$E_2$};
        
        % Beziehung
        \node[relation] ($R$) at (3,0) {$R$};
        
        % Kardinalitäten
        \node at (1.3,0.2) {\small1};
        \node at (4.7,0.2) {\small$N$};
        
        % Kanten
        \draw (prof) -- (lesen);
        \draw (lesen) -- (vl);
        
    \end{tikzpicture}
    \vspace{0.5cm}
    \begin{itemize}    
        \item Jeder Entität aus $E_1$ sind beliebig viele $E_2$ zugeordnet.
        \item Jeder Entität aus $E_2$ ist höchstens ein $E_1$ zugeordnet.
    \end{itemize}
    
    
    
    
\end{frame}

%%%%%%%%%%%%%%%%%%%%%%%%%%%%%%%%%%%%%%%%%%%%%%%%%%%%%%%%%%%%
% Kardinalitäten: N:M
\begin{frame}{Kardinalitäten: $N$ zu $M$}
    
    \centering
    \begin{tikzpicture}    
        
        % Entitäten
        \node[entity] (prof) at (0,0) {$E_1$};
        \node[entity] (vl) at (6,0) {$E_2$};
        
        % Beziehung
        \node[relation] ($R$) at (3,0) {$R$};
        
        % Kardinalitäten
        \node at (1.3,0.2) {\small$M$};
        \node at (4.7,0.2) {\small$N$};
        
        % Kanten
        \draw (prof) -- (lesen);
        \draw (lesen) -- (vl);
        
    \end{tikzpicture}
    
    \vspace{0.5cm}
    
    \begin{itemize}    
        \item Jeder Entität aus $E_1$ sind beliebig viele $E_2$ zugeordnet.
        \item Jeder Entität aus $E_2$ sind beliebig viele $E_1$ zugeordnet.
    \end{itemize}
    
    
    
    
\end{frame}


%%%%%%%%%%%%%%%%%%%%%%%%%%%%%%%%%%%%%%%%%%%%%%%%%%%%
% Beispiel Kardinalitäten
\begin{frame}{Beispiel Kardinalitäten}
    \includegraphics[width=\linewidth]{\BILDERPFAD/ERD_Uni.png}
\end{frame}

%%%%%%%%%%%%%%%%%%%%%%%%%%%%%%%%%%%%%%%%%%%%%%%%%%%%
% Beispiel Kardinalitäten
\begin{frame}{Beispiel Kardinalitäten}
    \includegraphics[width=\linewidth]{\BILDERPFAD/ERD_Uni_Kardinalitäten.png}
\end{frame}

%%%%%%%%%%%%%%%%%%%%%%%%%%%%%%%%%%%%%%%%%%%%%%%%%%%%
% (min,max) Notation
\begin{frame}{(min,max) Notation}
    
    \vspace{1cm}
    \centering
    \begin{tikzpicture}    
        
        % Entitäten
        \node[entity] (prof) at (0,0) {Professoren};
        \node[entity] (vl) at (6,0) {Vorlesung};
        
        % Beziehung
        \node[relation] (lesen) at (3,0) {lesen};
        
        % Kardinalitäten
        \node at (1.3,0.2) {\small$1$};
        \node at (4.7,0.2) {\small$N$};
        
        % Kanten
        \draw (prof) -- (lesen);
        \draw (lesen) -- (vl);
        
    \end{tikzpicture}
    
    \begin{tikzpicture}    
        
        % Entitäten
        \node[entity] (prof) at (0,0) {Professoren};
        \node[entity] (vl) at (6,0) {Vorlesung};
        
        % Beziehung
        \node[relation] (lesen) at (3,0) {lesen};
        
        % Kardinalitäten
        \node at (1.3,0.2) {\small$1$};
        \node at (4.7,0.2) {\small$N$};
        
        \node at (1.5,-0.2) {(0,*)};
        \node at (4.5,-0.2) {(1,1)};
        
        % Kanten
        \draw (prof) -- (lesen);
        \draw (lesen) -- (vl);
        
    \end{tikzpicture}
\end{frame}



%%%%%%%%%%%%%%%%%%%%%%%%%%%%%%%%%%%%%%%%%%%%%%%%%%%%
% Beispiel (min,max) Notation
\begin{frame}{(min,max) Notation}
    \includegraphics[width=\linewidth]{\BILDERPFAD/ERD_Uni_Kardinalitäten.png}
\end{frame}

